% p145_1; source PDF page 145; manually verified editable table.
\begin{tikzpicture}[x=1cm,y=-1cm,inner sep=0pt,outer sep=0pt]
\path[use as bounding box] (0,0) rectangle (16.000,8.335);
\draw[line width=.4pt] (0.000,0.000) rectangle (2.200,0.855);
\node[align=center,text width=1.940cm,font={\linespread{1}\fontsize{10}{13.50}\selectfont\bfseries}] at (1.100,0.427) {案情};
\draw[line width=.4pt] (2.200,0.000) rectangle (6.600,0.855);
\node[align=center,text width=4.140cm,font={\linespread{1}\fontsize{10}{13.50}\selectfont\bfseries}] at (4.400,0.427) {正犯的表现};
\draw[line width=.4pt] (6.600,0.000) rectangle (16.000,0.855);
\node[align=center,text width=9.140cm,font={\linespread{1}\fontsize{10}{13.50}\selectfont\bfseries}] at (11.300,0.427) {共犯的结论};
\draw[line width=.4pt] (0.000,0.855) rectangle (2.200,8.335);
\node[align=center,text width=1.940cm,font={\linespread{1}\fontsize{10}{13.50}\selectfont}] at (1.100,4.595) {甲教唆或帮助乙犯罪};
\draw[line width=.4pt] (2.200,0.855) rectangle (6.600,1.710);
\node[align=left,text width=4.140cm,font={\linespread{1}\fontsize{10}{13.50}\selectfont}] at (4.400,1.282) {乙根本未去犯罪，无罪};
\draw[line width=.4pt] (6.600,0.855) rectangle (16.000,1.710);
\node[align=center,text width=9.140cm,font={\linespread{1}\fontsize{10}{13.50}\selectfont}] at (11.300,1.282) {甲无罪};
\draw[line width=.4pt] (2.200,1.710) rectangle (6.600,3.115);
\node[align=left,text width=4.140cm,font={\linespread{1}\fontsize{10}{13.50}\selectfont}] at (4.400,2.413) {乙在预备阶段，犯罪中止};
\draw[line width=.4pt] (6.600,1.710) rectangle (16.000,3.115);
\node[align=left,text width=9.140cm,font={\linespread{1}\fontsize{10}{13.50}\selectfont}] at (11.300,2.413) {（1）甲构成犯罪预备。（2）甲若有中止行为且消除贡献，则构成犯罪中止};
\draw[line width=.4pt] (2.200,3.115) rectangle (6.600,3.970);
\node[align=left,text width=4.140cm,font={\linespread{1}\fontsize{10}{13.50}\selectfont}] at (4.400,3.543) {乙在预备阶段，犯罪预备};
\draw[line width=.4pt] (6.600,3.115) rectangle (16.000,3.970);
\node[align=center,text width=9.140cm,font={\linespread{1}\fontsize{10}{13.50}\selectfont}] at (11.300,3.543) {甲构成犯罪预备};
\draw[line width=.4pt] (2.200,3.970) rectangle (6.600,5.375);
\node[align=left,text width=4.140cm,font={\linespread{1}\fontsize{10}{13.50}\selectfont}] at (4.400,4.673) {乙在实行阶段，犯罪中止};
\draw[line width=.4pt] (6.600,3.970) rectangle (16.000,5.375);
\node[align=left,text width=9.140cm,font={\linespread{1}\fontsize{10}{13.50}\selectfont}] at (11.300,4.673) {（1）甲构成犯罪未遂。（2）甲若有中止行为且消除贡献，则构成犯罪中止};
\draw[line width=.4pt] (2.200,5.375) rectangle (6.600,6.230);
\node[align=left,text width=4.140cm,font={\linespread{1}\fontsize{10}{13.50}\selectfont}] at (4.400,5.803) {乙在实行阶段，犯罪未遂};
\draw[line width=.4pt] (6.600,5.375) rectangle (16.000,6.230);
\node[align=center,text width=9.140cm,font={\linespread{1}\fontsize{10}{13.50}\selectfont}] at (11.300,5.803) {甲构成犯罪未遂};
\draw[line width=.4pt] (2.200,6.230) rectangle (6.600,8.335);
\node[align=left,text width=4.140cm,font={\linespread{1}\fontsize{10}{13.50}\selectfont}] at (4.400,7.283) {乙在实行阶段，犯罪既遂};
\draw[line width=.4pt] (6.600,6.230) rectangle (16.000,8.335);
\node[align=left,text width=9.140cm,font={\linespread{1}\fontsize{10}{13.50}\selectfont}] at (11.300,7.283) {（1）甲对既遂结果若有贡献，则构成犯罪既遂。（2）甲若有中止行为且消除贡献，则构成犯罪中止，对乙的既遂结果不用负责};
\end{tikzpicture}
